\section{Related work}
\label{sec:related}

\paragraph{Joint Embedding Predictive Architecture (JEPA)}
\citet{LeCun2022} named the JEPA architecture, gave its two-encoder latent-prediction
diagram, and laid out the collapse taxonomy later studies inherit --- a self-supervised
encoder trained without care can converge to trivial solutions, so anti-collapse
mechanisms are a live engineering concern. Since then the field has moved toward using
JEPA as a world model, e.g.\ \citet{Maes2026LeWM}'s LeWorldModel.

\paragraph{Non-OLS methods for analysing regression structure.}
The pencil $(\Syx,\Sxx)$ this paper recovers is a classical object of study. \citet{Hotelling1936} introduced canonical correlation
analysis on exactly this generalized eigenproblem. Closer to
a dynamics-based claim, \citet{Varre2023} characterize the $t\to\infty$ limit of
gradient flow on a plain two-layer linear network as a min-norm interpolator --- an
ordering/timing result with no signed value and depth=$2$. Two recent JEPA-theoretic results also
recover structure without training dynamics: \citet{Klindt2026} prove identifiability
of LeJEPA's learned representation at the population optimum up to an orthogonal
transformation, while \citet{Cui2026} bound downstream planning regret via a low-rank spectral factorization
of an action-conditioned co-occurrence matrix, only once the encoder and predictor
already sit at their population-optimal solution. None of the above use training
dynamics as the recovery mechanism: the classical estimators require the full sample
covariance (or co-occurrence) structure in hand, and \citeauthor{Klindt2026}/
\citeauthor{Cui2026} characterize representations only once optimization has already
finished. The present paper recovers the same kind of object from training dynamics
alone, without ever forming $\hat\Sxx, \hat\Syx$ directly.

\paragraph{Analysing regression structure via model training dynamics.}
A separate line of work asks whether training dynamics themselves reveal this
structure, without a closed-form estimator. \citet{Littwin2024} shows gradient-flow
dynamics of a linear JEPA pair, with no added regularizer, preferentially learn
high-regression-coefficient features faster, at critical time
$\tilde t_r^* \asymp 1/(\lamstar \eps^{(2L-1)/L})$ — dynamics do encode regression
structure, at a rate in the same $\eps^{1/L}$-family as this paper's own
$O(\eps^{1/L}|\log\eps|)$ recovery rate (\S\ref{sec:estimating-spectrum}). But the
result assumes $\Sxx$ and $\Syx$ share an eigenbasis (known in advance) and recovers a
magnitude-linked critical time, not a signed coefficient - an assumption relaxed by \citet{Goh2026ordering}.
\citet{Domine2024} generalizes along a different axis — exact solutions across a continuous
$\lambda$-balanced initialization spectrum, not only the zero-balanced case — but
$\tilde\Sxx,\tilde\Syx$ stay fixed and known throughout. 
Together these three confirm training dynamics genuinely encode
regression structure but stop short on different axes. The present paper is the first to close all three at once,
recovering an unknown, signed structure from dynamics alone.

\paragraph{Solving training dynamics.}
This paper's own proof strategy inherits its central diagonalization step from
\citet{Saxe2014,Saxe2019}, who first showed that under aligned, orthogonal
initialization a deep linear network's coupled gradient-flow ODEs decouple into
per-mode scalar Bernoulli-type equations with closed-form sigmoidal trajectories; we
diagonalize onto the generalized eigenbasis instead of an ordinary one, yielding the
same Saxe-style depth exponents, but the decoupling step itself is classical and not
re-derived here. \citet{Arora2018,Arora2019} formalize the ``balanced initialization''
condition this paper's own Assumption~\ref{ass:aligned-init} restricts to, showing deep
linear gradient descent is equivalent to a preconditioned single-layer system with a
rigorous linear-rate convergence guarantee under approximately balanced initialization
and whitened data --- the optimization/convergence-rate side of the same initialization
concept \citeauthor{Saxe2014} solve exactly. \citet{Nam2025} instead argues,
independent of any new dynamics result, that layerwise-linear networks are the right
minimal lens for training-dynamics phenomena generally, directly supporting this
paper's own reason for studying the linear case first (\S\ref{ssec:positioning}).
Together these establish that deep linear training dynamics are exactly solvable and worth studying on purpose
even though OLS could already recover the population structure directly; none ask this
paper's own question of that same machinery: identifiability of an unknown, signed
structure.

\paragraph{Formal verification of learning theory.}
Machine-checked learning theory is a small but growing body of work.
\citet{Bagnall2019} started this line, formalizing PAC bounds in Coq for finite
hypothesis classes. \citet{Sonoda2026} made the major advance to infinite hypothesis
classes, in Lean~4 against Mathlib --- the same proof assistant this paper uses. A companion repository
(\url{https://github.com/davidcagoh/jepa-rho-recovery}, public) contains the Lean
development and the \textsc{PlateauRecover} reference implementation, alongside, not
in place of, the paper's main mathematical contribution. 